\documentclass[10pt,a4paper]{amsart}
\usepackage[margin=19mm]{geometry}
\usepackage{amsmath,amssymb,mathtools}
\usepackage[scr=rsfs,cal=euler]{mathalfa}
\usepackage{xfrac}
\usepackage[unicode,hidelinks,bookmarks=false]{hyperref}
\newcommand{\ie}{i.e.\ }
\newcommand{\cf}{cf.\ }
\newcommand{\resp}{respectively\ }
\newcommand{\Subsection}{\subsection}
\providecommand{\nobreakdash}{\nobreak\hbox{-}}
\setlength{\emergencystretch}{2em}
\multlinegap0pt
\usepackage[shortlabels]{enumitem}
\usepackage{bm}
\usepackage{amssymb}
\usepackage{mathtools}
\usepackage{tikz}
\newtheorem{thm}{Theorem}
\newtheorem{rmq}{Remark}
\newtheorem{prop}{Proposition}
\usepackage{cleveref}
\crefname{thm}{Theorem}{Theorems}
\crefname{rmq}{Remark}{Remarks}
\crefname{prop}{Proposition}{Propositions}
\newcommand{\Gammaplus}[1]{ {  \Gamma}^{(+)}_{#1}}
\newcommand{\Top}{\mathcal{T}}
\newcommand{\tH}{\widetilde{H}^{(q,t)}}
\title{Macdonald characters from a new formula for Macdonald polynomials}
\author{Houcine Ben Dali, Michele D'Adderio}
\date{Source: arXiv:2404.03904v2 (CC BY 4.0)}
\begin{document}
\maketitle

\noindent\emph{Source and licence.} Macdonald characters from a new formula for Macdonald polynomials. Version arXiv:2404.03904v2 (CC BY 4.0), distributed by arXiv under the Creative Commons Attribution 4.0 International licence, http://creativecommons.org/licenses/by/4.0/, as recorded in the arXiv licence metadata for this exact version and shown on the abstract page https://arxiv.org/abs/2404.03904v2. Full licence notice: \texttt{licenses/arxiv-2404.03904-CC-BY-4.0.txt}.\par\smallskip

\noindent\textit{Editorial scope.} Counted unit: the article's thm modified (source0.tex lines 408--415). The article's complete original proof of it is delivered with it (source0.tex lines 453--487, one \texttt{proof} environment of 35 lines, which the article places away from the statement). No mathematics is written, repaired or re-derived: the statement and the proof are the article's own lines, in the article's own notation, and no retained line is substituted or re-flowed.  Retained uncounted as the source-local context the proof needs: lines 278--280; lines 378--380; lines 425--427; lines 431--433; lines 436--450.  Nothing was omitted from any retained span, so the package carries no omission record.  No retained line contains a citation, so the wrapper carries no bibliography and no bibliography entry.  No retained line contains a figure environment, an image-inclusion command or drawing code, so no image is delivered and the package declares no asset dependency.  Editorial formatting only, in addition: an AMS \texttt{amsart} preamble that restates the article's own declarations and macros used by the retained lines, namely the environments \texttt{thm}, \texttt{rmq}, \texttt{prop} on separate counters, together with the standard packages those lines need.  The retained environments print in this document as Theorem 1, Rmq 1, Proposition 1 in order of appearance, whereas the article prints the same items with the numbering of its own full document.  The complete native source of the article as distributed in the arXiv e-print for this exact version is bundled unchanged as \texttt{sources/157253/source0.tex}.
\begin{equation}\label{eq n}
  n(\lambda):=\sum_{1\leq i\leq \ell(\lambda)}\lambda_{i}(i-1)=\sum_{\Box\in\lambda}\ell_\lambda'(\Box).  
\end{equation}

	\begin{align}\label{eq creation modified 2}
\Gamma(u,q^{\lambda_1})\Gamma(tu,q^{\lambda_2})\cdots \Gamma(t^{k-1}u,q^{\lambda_{k}})\cdot 1 & =\nabla \Top_{\frac{1}{u}} \widetilde{H}_\lambda[uX]=\nabla \widetilde{H}_\lambda[uX+1].
	\end{align}

\begin{equation}\label{eq Gamma formula}
  \Gamma(u,v)\cdot \tH_\lambda=\sum_{k\geq 0}u^k\sum_{\xi\vdash k+|\lambda|}\widetilde\eta_{\lambda,\xi}\tH_\xi\prod_{\Box\in\xi/ \lambda}\left(v-q^{a'(\Box)}t^{\ell'(\Box)}\right),  
\end{equation}

\begin{rmq}\label{rmq Box0}
    Let $\Box_0$ be the cell of coordinates $(1,s+1)$. This cell is characterized by $a'(\Box_0)=s$ and $\ell'(\Box_0)=0$. Notice that the condition $\lambda_1\leq s$ is equivalent to saying that $\Box_0$ is not a cell of the Young diagram of $\lambda$.
\end{rmq}

\begin{prop}\label{prop stability and vhanishing condition}
    Fix $s,k\geq 0$ two integers. 
    \begin{enumerate}[label={\normalfont(}\arabic*\normalfont)]
        \item The space $\Lambda_{\leq s}$ is stabilized by the action of the operator $[u^k]\Gamma(u,q^s)$.
        \item If $k>s$, then $[u^k]\Gamma(u,q^s)=0$ as operators on $\Lambda_{\leq s}$. 
    \end{enumerate} 
\begin{proof}
Fix a partition $\lambda$ such that $\lambda_1\leq s$.
From \cref{eq Gamma formula}, we know that 
\begin{equation}\label{eq prop stability}
  [u^k]\Gamma(u,q^s)\cdot \tH_\lambda=\sum_{\xi\vdash k+|\lambda|}\widetilde\eta_{\lambda,\xi}\tH_\xi\prod_{\Box\in\xi/ \lambda}\left(q^s-q^{a'(\Box)}t^{\ell'(\Box)}\right),  
\end{equation}
where the sum is taken over horizontal strips $\xi/\lambda$ of size $k$. Notice that, with the notation of \cref{rmq Box0}, the quantity $q^s-q^{a'(\Box)}t^{\ell'(\Box)}$ is 0 if and only if $\Box=\Box_0$. Moreover, if $\xi$ is a partition such that $\xi_1>s$, then $\Box_0\in\xi/\lambda$ and therefore, $\xi$ does not contribute to the sum in \cref{eq prop stability}. This proves $(1)$. If $k>s$ then any horizontal strip $\xi/\lambda$ of size $k>s$ contains necessarily the cell $\Box_0$. This gives (2).    
\end{proof}
\end{prop}

\begin{thm} \label{thm creation modified}
	For any partition $\lambda=[\lambda_1,\lambda_2,\dots,\lambda_k]$, we have
	\begin{equation}\label{eq thm creation modified}
		\Gammaplus{\lambda_1}\Gammaplus{\lambda_2}\cdots\Gammaplus{\lambda_k}\cdot1=t^{-n(\lambda)}\tH_{\lambda}
  \end{equation}
  \noindent where
  $$\Gammaplus{m}:=[u^m]\nabla^{-1}\Gamma(u,q^m)\nabla.$$
\end{thm}

\begin{proof}[Proof of \cref{thm creation modified}]
    Let $m$ and $\ell$ denote respectively the size and the length of $\lambda$. Extracting the coefficient of $u^m$ in  \cref{eq creation modified 2}, we get 
\begin{align}
	\sum_{m_1+\dots+m_\ell=m}\left([u^{m_1}]\Gamma(u,q^{\lambda_1})\right)\cdots \left([u^{m_\ell}]\Gamma(t^{\ell-1}u,q^{\lambda_{\ell}}))\right)\cdot 1 & =\nabla \tH_\lambda[X].
\end{align}
Hence,
\begin{align}\label{eq formula modified}	\sum_{m_1+\dots+m_\ell=m}\left([u^{m_1}]\nabla^{-1}\Gamma(u,q^{\lambda_1})\nabla\right)\cdots \left([u^{m_\ell}]\nabla^{-1}\Gamma(t^{\ell-1}u,q^{\lambda_{\ell}})\nabla\right)\cdot 1 & =\tH_\lambda[X].
\end{align}
We start by proving that the only tuple $(m_1,\dots,m_\ell)$ which contributes to this sum is $(\lambda_1,\dots,\lambda_\ell).$
First, notice that for each tuple $(m_1,\dots,m_{\ell})$, we obtain by induction and using item (1) of \cref{prop stability and vhanishing condition} that for any $0\leq i \leq \ell$ we have
\begin{equation}\label{eq stability induction}
  \left([u^{m_{\ell-i}}]\nabla^{-1}\Gamma(t^{\ell-i-1}u,q^{\lambda_{\ell-i}})\nabla\right)\cdots \left([u^{m_\ell}]\nabla^{-1}\Gamma(t^{\ell-1}u,q^{\lambda_{\ell}})\nabla\right)\cdot 1\in \Lambda_{\leq \lambda_{\ell-i}}.  
\end{equation}
Moreover, if for some $1\leq i\leq \ell$, we have $m_i>\lambda_i$, then from \cref{prop stability and vhanishing condition} item (2) and \cref{eq stability induction} we have  
\begin{equation}
  \left([u^{m_{i}}]\nabla^{-1}\Gamma(t^{i-1}u,q^{\lambda_{i}})\nabla\right)\cdots \left([u^{m_\ell}]\nabla^{-1}\Gamma(t^{\ell-1}u,q^{\lambda_{\ell}})\nabla\right)\cdot 1=0.  
\end{equation}
We deduce that any tuple $(m_1,\dots,m_\ell)$ which is different from $(\lambda_1,\dots,\lambda_\ell)$ does not contribute to the sum of \cref{eq formula modified}. 
Therefore
\begin{align*}
\left([u^{\lambda_1}]\nabla^{-1}\Gamma(u,q^{\lambda_1})\nabla\right)
\left([u^{\lambda_2}]\nabla^{-1}\Gamma(tu,q^{\lambda_2})\nabla\right)\cdots \left([u^{\lambda_\ell}]\nabla^{-1}\Gamma(t^{\ell-1}u,q^{\lambda_{\ell}})\nabla\right)\cdot 1 & =\tH_\lambda[X].
\end{align*}
This can be rewritten as follows,
\begin{multline*}
\left([u^{\lambda_1}]\nabla^{-1}\Gamma(u,q^{\lambda_1})\nabla\right)
\left(t^{\lambda_2}[u^{\lambda_2}]\nabla^{-1}\Gamma(u,q^{\lambda_2})\nabla\right)\cdots \\
\left(t^{\lambda_\ell(\ell-1)}[u^{\lambda_\ell}]\nabla^{-1}\Gamma(u,q^{\lambda_{\ell}})\nabla\right)\cdot 1 =\tH_\lambda[X].
\end{multline*}
Using the definition of $n(\lambda)$ (see \cref{eq n}), we deduce that
\begin{align*}
t^{n(\lambda)}\left([u^{\lambda_1}]\nabla^{-1}\Gamma(u,q^{\lambda_1})\nabla\right)\cdots \left([u^{\lambda_\ell}]\nabla^{-1}\Gamma(u,q^{\lambda_{\ell}})\nabla\right)\cdot 1 & =\tH_\lambda[X].
\end{align*}
This completes the proof of the theorem.
\end{proof}

\end{document}
